% TOC entry added in LTR context → title left, page number right (like all other entries)
\phantomsection
\addcontentsline{toc}{chapter}{\textarabic{ملخص}}

% Page content in RTL
\begin{Arabic}

\chapter*{ملخص}
\markboth{\textarabic{ملخص}}{\textarabic{ملخص}}

تعتمد أنظمة المعلومات القرارية الكلاسيكية على معماريات مستودعات البيانات المصممة لبيئات
أعمال مستقرة واحتياجات تحليلية ثابتة. في مواجهة تسارع التحولات التنظيمية، وتكاثر مصادر
البيانات غير المتجانسة، والنضج التحليلي المتزايد للمنظمات، تصبح هذه المعماريات الصارمة
عائقاً أمام تطور أنظمة دعم القرار.

تعالج هذه الأطروحة مشكلة قابلية التوسع في معماريات اتخاذ القرار على محورين متكاملين:
\textbf{قابلية التوسع الهيكلية}، التي تتعلق بقدرة النظام على دمج مصادر جديدة وإدارة
تغييرات المخطط والتكيف مع نمو الحجم دون تعطيل المعمارية القائمة؛ و\textbf{قابلية التوسع
التحليلية}، التي تعكس التقدم من ذكاء الأعمال الوصفي البحت نحو القدرات التنبؤية والتوصيفية
المدمجة بشكل أصيل في خط أنابيب القرار.

لمعالجة هذه التحديات، يعتمد هذا البحث على منهجية Data Vault~2.0 كأساس لطبقة النمذجة
الرشيقة، ويستكشف أنماط دمج النماذج التحليلية المتقدمة ضمن المعماريات الحديثة (Lakehouse،
Delta Architecture). تقدم الأطروحة ثلاثة إسهامات مفاهيمية: إطار عمل رسمي للتوسعية
الثنائية (C1)، وتحليل نقدي لإمكانية تركيب Data Vault~2.0 مع نموذج Lakehouse مع أداة
معايير (C2)، ومعمارية مرجعية متعددة الطبقات مع عقد تكامل رسمي يُقيَّم وفق ستة معايير (C3).

\bigskip

\noindent\textbf{الكلمات المفتاحية:}
معمارية قابلة للتوسع، نظام معلومات قراري، Data Vault~2.0، مستودع بيانات رشيق،
تحليلات تنبؤية، ذكاء الأعمال، Lakehouse، ETL/ELT، تطور المخطط، تعلم الآلة.

\end{Arabic}
